\subsection{RQ1: Does sequential PPO outperform contextual learning?}
\label{sec:rq1}

Table~\ref{tab:main} reports all policies on the 140 test episodes. The bandit
reaches a mean return of \TestBandit{} [\TestBanditLo, \TestBanditHi] over
\TestBanditRoots{} training roots and PPO \TestPPO{} [\TestPPOLo, \TestPPOHi]
over \TestPPORoots{}. On paired episodes the bandit's advantage is \Gap{}
($t$-interval over roots [\GapTLo, \GapTHi]; two-stage bootstrap
[\GapLo, \GapHi]); it is positive on \GapRootsPositive{} of \GapRoots{} roots
and in \GapWinRate\,\% of paired episodes (paired $d_z=\GapDz$, pooled over roots and episodes). The ordering does not
depend on checkpoint selection: with the final 400k checkpoints the gap is
\GapFinal{} [\GapFinalLo, \GapFinalHi] (\GapFinalRootsPositive{} of \GapRoots{}
roots).
Per scenario (Table~\ref{tab:perscen}), the bandit is ahead in every
scenario except the double-failure scenario, where the difference is small
and its interval includes zero (\GapDoublefailure{}); the
largest gap is on the 24-hour \texttt{full\_day} scenario. Both learners
dominate the static, random and CSPF baselines. Against greedy the bandit is
ahead on \BanditMinusGreedyRootsPositive{} of \GapRoots{} roots
(\BanditMinusGreedy{} [\BanditMinusGreedyLo, \BanditMinusGreedyHi]), PPO on
\PPOMinusGreedyRootsPositive{} of \GapRoots{} (\PPOMinusGreedy{}
[\PPOMinusGreedyLo, \PPOMinusGreedyHi]).

We also compare with a controller that does not learn at all. MILP-track, which solves a min-max-utilization path-selection problem on
current telemetry every interval and moves one demand towards the solution,
reaches \TestMilp{} [\TestMilpLo, \TestMilpHi]. We detect no difference in return
between the bandit and MILP-track at five roots (bandit $-$ MILP-track
$=\BanditMinusMilp{}$, $t$-interval [\BanditMinusMilpTLo,
\BanditMinusMilpTHi], positive on \BanditMinusMilpRootsPositive{} of
\GapRoots{} roots); this is not an equivalence claim, and the two differ by
scenario (Table~\ref{tab:milp_perscen}): on the double-failure
scenario the bandit's difference is \BanditMilpDoublefailure{} and positive on
\BanditMilpDoublefailureRootsPositive{} of \GapRoots{} roots. PPO falls
short of MILP-track ($\PPOMinusMilp$, $t$-interval [\PPOMinusMilpTLo,
\PPOMinusMilpTHi], \PPOMinusMilpRootsPositive{} of \GapRoots{} roots
positive). Operationally, the
bandit delivers \TestBanditDelivered\,\% of offered traffic with
\TestBanditSla{} SLA-violating demand-intervals per episode and
\TestBanditReroutes{} reroutes per hour, against \TestPPODelivered\,\%,
\TestPPOSla{} and \TestPPOReroutes{} for PPO and \TestMilpDelivered\,\%,
\TestMilpSla{} and \TestMilpReroutes{} for MILP-track; the return gap to PPO
thus corresponds to about half a percentage point of delivered traffic, fewer
SLA violations and a third of PPO's churn. MILP-track attains higher network utility and pays more in
move costs, which its objective does not contain (Table~\ref{tab:decomp}); the
bandit reaches a similar return with \ChurnRatioBanditMilp{} times its reroute
rate. MILP-track's two parameters (the minimum utilization gain that triggers a
move, and which mismatched demand to move) were not tuned for this reward.
Selecting them on the validation seeds (eight combinations; minimum gain
\MilpTunedGain{}, largest-volume demand first) gives a controller that reaches
\TestMilpTuned{} [\TestMilpTunedLo, \TestMilpTunedHi] on the test episodes
with \TestMilpTunedReroutes{} reroutes per hour, ahead of the bandit by
\MilpTunedAhead{} ($t$-interval [\MilpTunedAheadTLo, \MilpTunedAheadTHi];
the bandit is ahead on \BanditMinusMilpTunedRootsPositive{} of \GapRoots{}
roots) and of PPO by more. With a model of the objective and a little tuning,
per-interval optimization beats both learners here.

\begin{table*}[t]
\centering\small
\caption{Test results (7 scenarios $\times$ 20 seeds, 140 paired episodes).
Learners: mean over training roots of the selected checkpoint, two-stage
bootstrap 95\,\% CI; fixed policies: scenario-stratified bootstrap CI.
Per-interval = return divided by the number of decisions, averaged over
episodes. SLA = SLA-violating demand-intervals per episode. $\dagger$
clairvoyant reference (sees future traffic), not a controller.}
\label{tab:main}
\begin{adjustbox}{max width=\textwidth}
\begin{tabular}{@{}lrrrrrrr@{}}
\toprule
Policy & Return [95\,\% CI] & Per-interval & Delivered & SLA & Mean max-util & Reroutes/h & Moved Mb/s \\
\midrule
\tabinput{tables/main_results}
\bottomrule
\end{tabular}
\end{adjustbox}
\end{table*}

\begin{table}[t]
\centering\small
\caption{Bandit $-$ PPO paired return difference per scenario (selected
checkpoints; two-stage bootstrap 95\,\% CI over roots, then episodes), number
of roots with a positive difference, and share of paired episodes the bandit
wins.}
\label{tab:perscen}
\begin{tabular}{@{}lrrr@{}}
\toprule
Scenario & Difference [95\,\% CI] & Roots $+$ & Episodes won \\
\midrule
\tabinput{tables/main_per_scenario}
\bottomrule
\end{tabular}
\end{table}
